\subsection{Solution Procedures}
\label{sec:station-procedures}

\subsubsection{Active power given at the PCC}

With $P_{pcc}$ given, the AC network, the AC line, the transformer and the converter are calculated once in this
order. The converter needs its DC voltage, which depends on the DC current through the DC reactor
(Section~\ref{sec:station-dc}), and the DC current follows from the converter. This loop is solved by a fixed-point
iteration on $U_{dpe\_conv}$:
\begin{enumerate}
  \item start with $U_{dpe\_conv} = U_{dpe}$;
  \item calculate the chain with $U_{dc} = k \cdot U_{dpe\_conv}$, giving $I_d$;
  \item update $U_{dpe\_conv} = U_{dpe} \pm R_{dcr} \cdot I_d$ and repeat from step 2 until the change is below
        $10^{-10}$\,kV.
\end{enumerate}
Without a DC reactor one pass is enough. The losses of the converter depend on its DC current
$I_{dc} = (P_{ac} + P_{loss\_cv}) / U_{dc}$; inside the converter they are found by a fixed-point iteration as
well (Section~\ref{sec:converter-losses}).

\subsubsection{Active power given at the DC terminal}

With $P_{dc}$ given at the station DC terminal (all poles), the DC side is known directly:
\begin{equation}
  I_d = \frac{|P_{dc}|}{n_{conv} \cdot k \cdot U_{dpe}} , \qquad
  U_{dpe\_conv} = U_{dpe} \pm R_{dcr} \cdot I_d , \qquad
  P_{dc\_cv}^{target} = \mp k \cdot U_{dpe\_conv} \cdot I_d ,
\end{equation}
upper signs for a rectifier ($P_{dc} < 0$). The PCC power is the root of
\begin{equation}
  F(P_{pcc}) = P_{dc\_cv}(P_{pcc}) - P_{dc\_cv}^{target} ,
\end{equation}
where $P_{dc\_cv}(P_{pcc})$ is the converter DC power from the AC chain at the fixed converter voltage. The root is
found by the secant method, starting from $P_{pcc} = n_{conv} \cdot P_{dc\_cv}^{target}$ and 1.01 times that value,
with the tolerance $10^{-9}$\,MW and at most 50 iterations.

\subsubsection{Station totals and checks}

The station losses are
\begin{equation}
  P_{loss\_st} = n_{conv} \cdot \left( P_{loss\_tr} + P_{aux}/n_{conv} + P_{loss\_cv} \right)
                 + 2 \cdot R_{dcr} \cdot I_d^2 + P_{loss\_line} ,
\end{equation}
and the efficiency is the output power divided by the input power: $|P_{dc}| / |P_{pcc}|$ for a rectifier,
$P_{pcc} / P_{dc}$ for an inverter. The result carries the operation that follows from the sign of the given power
and warnings for a PCC voltage outside its band, an AC line or transformer loading above 100\,\% and every exceeded
limit of the converter. A rectifier whose converter would not deliver DC power, or the reverse for an inverter, is
rejected.
